\section{Implementation notes (aligned to \citrees)}
\label{app:implementation}

This appendix records implementation details that matter for reproducibility and
for interpreting the scope of the theoretical statements in
\cref{sec:theory}.

This appendix complements Appendix~\ref{app:methods}: the methods appendix
records algorithmic details, while this appendix focuses on reproducibility
details (RNG, compilation modes, and implementation choices that affect how to
interpret results).

\subsection{Registry pattern}

\citrees uses a registry mechanism that maps short identifiers for association
statistics and split criteria to their corresponding implementations and
permutation-test adapters. This keeps the codebase extensible while preserving a
small, stable external interface.

\subsection{Testing modes (compiled/uncompiled)}

Performance-critical routines are compiled for speed. The test suite supports
both compiled execution (to validate optimized code paths) and an uncompiled
execution mode (to enable line coverage measurement in CI).

\subsection{Randomness and reproducibility}

\citrees supports reproducible randomness via a user-provided base seed. In
parallel permutation testing, the implementation uses a per-resample seeding
pattern inside parallel compiled loops. In particular, it uses
%
\begin{equation}
  s_i := s_0 + i
\end{equation}
%
to ensure deterministic per-resample pseudo-randomness across parallel workers.
This convention is an implementation constraint: generator-based RNG streams are
not always supported in all parallel compiled contexts.

\subsection{Bootstrap methods in forests}

Forests can optionally resample training rows per tree. In bootstrap mode,
rows are sampled with replacement uniformly. This affects only the per-tree
resampling distribution, and it can invalidate the root-level exchangeability
assumption A0.2 used in \cref{sec:theory} because duplicated indices induce
non-permutation-invariant equality constraints in the bootstrap sample.
Accordingly, the fixed-node/root validity statements apply only to trees or
nodes for which A0.2 still holds after resampling (for example, full-data roots
or subsampling without replacement).
